\documentclass[11pt]{article}

\usepackage[margin=1in]{geometry}
\usepackage{amsmath,amssymb}
\usepackage{hyperref}
\usepackage{enumitem}

\hypersetup{
  colorlinks=true,
  linkcolor=blue,
  urlcolor=blue,
  citecolor=blue
}
\emergencystretch=2em

\newcommand{\domop}{\operatorname{dom}}
\newcommand{\inte}{\operatorname{int}}
\newcommand{\sta}{\operatorname{star}}
\newcommand{\gra}{\operatorname{gra}}
\newcommand{\To}{\rightrightarrows}

\title{A Literature-Implied Convex-Domain Answer to an FPV Sum Problem}
\author{Math discovery packet}
\date{June 28, 2026}

\begin{document}
\maketitle

\section{Status}

This is a lightweight \emph{literature-implied answer}, not an original
solution.  The source paper is J. M. Borwein and L. Yao,
\emph{Recent progress on Monotone Operator Theory}, arXiv:1210.3401.
In Section ``Open problems in Monotone Operator Theory'', p. 33 of the arXiv
PDF, the paper asks:
\begin{quote}
Let \(A,B:X\To X^*\) be maximally monotone with
\(\domop A\cap\inte\domop B\ne\varnothing\). Assume that \(A\) is of type
(FPV). Is \(A+B\) necessarily maximally monotone?
\end{quote}

The supporting paper is J. M. Borwein and L. Yao,
\emph{Sum theorems for maximally monotone operators of type (FPV)},
arXiv:1305.6691.

\section{Identification}

The supporting paper proves the following theorem and corollary.
\begin{enumerate}[leftmargin=*]
\item Theorem 4.2, the ``(FPV) Sum Theorem'', proves maximal monotonicity of
\(A+B\) when \(A\) and \(B\) are maximally monotone,
\(\sta(\domop A)\cap\inte\domop B\ne\varnothing\), and \(A\) is of type
(FPV).
\item Corollary 4.3, ``Convex domain'', states that if
\(\domop A\cap\inte\domop B\ne\varnothing\) and \(A\) is of type (FPV) with
convex domain, then \(A+B\) is maximally monotone.
\end{enumerate}

Since convexity of \(\domop A\) implies \(\sta(\domop A)=\domop A\), Corollary
4.3 gives a direct affirmative answer to the source question in the
convex-domain subcase.

\section{Scope}

The original arXiv:1210.3401 question does not assume that \(\domop A\) is
convex.  The supporting paper does not claim to settle that general case:
after Corollary 4.3 it explicitly asks whether the star-shaped hypothesis on
\(\domop A\) in Theorem 4.2 can be relaxed or dropped, and then re-poses the
unqualified FPV sum problem.

A later arXiv paper, arXiv:1411.6199, claimed to settle the unqualified
question, but the arXiv record states that the paper has been withdrawn by the
author because of an error in Lemma 1.  This packet therefore records only the
valid convex-domain literature subcase.

\section{References}

\begin{itemize}[leftmargin=*]
\item J. M. Borwein and L. Yao, \emph{Recent progress on Monotone Operator
Theory}, arXiv:1210.3401.
\item J. M. Borwein and L. Yao, \emph{Sum theorems for maximally monotone
operators of type (FPV)}, arXiv:1305.6691.
\item S. R. Pattanaik and D. K. Pradhan, \emph{Maximality of the sum of the
monotone operator of type (FPV) and a maximal monotone operator},
arXiv:1411.6199, withdrawn.
\end{itemize}

\end{document}
